\chapter{Specialists}
\noindent Select the model first; copy only required snippets and dependencies.
The contest prelude supplies \texttt{ll, vi, pii, rep, all, sz} and standard
headers. SCC and 2-SAT use GNU C++ syntax. Globals such as \texttt{dfs, Time,
P, mod} can collide: rename or namespace combined snippets. Header test status
is upstream evidence, not a proof of your reduction.

\begin{lstlisting}
#include <bits/stdc++.h>
using namespace std;
#define rep(i,a,b) for(int i=(a);i<(b);++i)
#define all(x) begin(x),end(x)
#define sz(x) int((x).size())
using ll=long long;
using vi=vector<int>;
using pii=pair<int,int>;
\end{lstlisting}

\section{Strings}
\subsection{One pattern and borders}
\texttt{match(text,pat)} returns zero-based occurrence starts, including overlaps;
use a nonempty pattern. Its nul separator must be absent from the inputs.
\texttt{pi(s)[i]} is the longest proper prefix ending at $i$.
A repeated block has candidate length $p=n-\texttt{pi(s).back()}$;
full repetition additionally requires $n\bmod p=0$. Time $O(n+m)$.
\sourceimport{strings}{KMP.h}

\subsection{Prefix matches at every suffix}
\texttt{Z(s)[i]} is the common-prefix length of $s$ and suffix $i$;
\texttt{z[0]=0}. Time/space $O(n)$. For pattern--separator--text matching,
choose a separator outside both alphabets.
\sourceimport{strings}{Zfunc.h}

\subsection{Many fixed patterns}
Construct \texttt{AhoCorasick ac(patterns)} from a mutable vector.
Default alphabet is uppercase A--Z: change \texttt{first} and \texttt{alpha}
before passing other symbols. Empty patterns are forbidden; duplicates allowed.
\texttt{find(text)} returns the longest ending pattern index or $-1$ at each
position. \texttt{findAll(patterns,text)} groups \emph{all} matches by starting
position. Enumerating matches costs their output size. After construction,
missing trie transitions have become failure-link transitions; do not treat
\texttt{next} as a sparse trie. Build $O(26L)$, scan $O(n)$.
\sourceimport{strings}{AhoCorasick.h}

\subsection{Other source choices}
\textbf{Hashing.h:} \texttt{HashInterval h(s); h.hashInterval(l,r)} hashes
$[l,r)$ in $O(1)$ after $O(n)$ preprocessing. Uses arithmetic modulo $2^{64}-1$;
collisions remain possible. Compare lengths; prefer KMP/Z for exact matching.
\textbf{Manacher.h:} \texttt{p=manacher(s)} gives even lengths
$2p[0][i]$ and odd lengths $2p[1][i]+1$.
\textbf{SuffixArray.h:} \texttt{SuffixArray a(s)} includes empty suffix
\texttt{a.sa[0]=n}; \texttt{a.lcp[i]} compares adjacent sorted suffixes,
not arbitrary text positions. No nul input; ordinary ASCII avoids signed-char
alphabet issues. Build $O(n\log n)$. These three sources are mapped, not printed here.

\section{Directed constraints}
\subsection{Strong components}
\texttt{scc(g,callback)} writes global \texttt{comp[v]} and \texttt{ncomps}.
Inter-component edges go from larger to smaller component IDs (reverse
topological order). Copy callback blocks if storing them: the vector is reused.
Build the condensation DAG after IDs exist. Recursive DFS: check stack capacity.
Time $O(V+E)$.
\sourceimport{graph}{SCC.h}

\subsection{Boolean pairwise clauses}
A literal is variable index $x$ or bitwise negation \texttt{\~{}x};
\texttt{!x} is wrong. \texttt{either(a,b)} adds $a\lor b$;
encode $a\Rightarrow b$ as \texttt{either(\~{}a,b)}.
\texttt{setValue(a)} forces literal $a$ true.
Read \texttt{values} only after \texttt{solve()} succeeds.
\texttt{atMostOne} supplies no at-least-one constraint. General clauses of
three or more literals do not fit 2-SAT. Time $O(V+E)$; its SCC code is internal.
\sourceimport{graph}{2sat.h}

\subsection{Undirected blocks: source contract}
\textbf{BiconnectedComponents.h:} global \texttt{ed[u]} stores
\texttt{(neighbor,edgeId)}. Insert an undirected edge twice with one shared ID;
assign different IDs to parallel physical edges. \texttt{bicomps(callback)}
reports edge blocks; vertices can belong to several blocks. Bridges are only
recognized at the marked branch: add collection code there. Handle self-loops
separately. Articulation vertices need additional root-child bookkeeping.
Time $O(V+E)$. Source is not printed in this appendix.

\section{Matching and capacities}
\subsection{Bipartite matching}
\texttt{g} contains neighbors of left vertices, indexed within the right side.
Initialize \texttt{vi r(rightSize,-1)}; \texttt{hopcroftKarp(g,r)} returns
cardinality and writes \texttt{r[right]=left}, or $-1$. Time $O(E\sqrt V)$.
This is unweighted bipartite matching, not general or weighted matching.
\sourceimport{graph}{HopcroftKarp.h}

\subsection{Flow and residual cut}
\texttt{Dinic f(n); f.addEdge(u,v,cap)} adds a directed edge with zero initial
reverse capacity. Require nonnegative capacities, $s\ne t$, and bounded total
flow fitting \texttt{ll}. Avoid self-loops unless adapting reverse indices.
\texttt{f.calc(s,t)} mutates residual capacities and returns \emph{additional}
flow. Afterwards \texttt{leftOfMinCut(v)} tests residual source reachability.
\texttt{Edge.c} is residual capacity, \texttt{oc} original, \texttt{flow()}
nonnegative net usage. Header bound $O(VE\log U)$; DFS is recursive.

Matching: source$\to$left, compatible left$\to$right, right$\to$sink,
all capacity one. Vertex capacity: split into in$\to$out. Infinite capacity
must exceed every finite cut with a proved safe bound; do not use
\texttt{LLONG\_MAX} blindly. Min-cost objectives require a different solver.
This header exists in source but is omitted from the original default graph
chapter; it is explicitly printed here.
\sourceimport{graph}{Dinic.h}

\section{Tree paths}
\subsection{Ancestor jumps}
Supply \texttt{parent[root]=root}, valid parent indices, and separate depths.
\texttt{up=treeJump(parent)}, \texttt{jmp(up,v,k)},
\texttt{lca(up,depth,u,v)}. Restrict jumps to supported bits/depth.
Build/storage $O(n\log n)$, query $O(\log n)$.
\sourceimport{graph}{BinaryLifting.h}

\subsection{Heavy-light source contract}
\textbf{HLD.h:} nonempty connected tree, root zero.
\texttt{HLD<false>} stores vertices; \texttt{HLD<true>} stores edges at their
deeper endpoints. \texttt{modifyPath(u,v,x)} adds $x$;
\texttt{queryPath(u,v)} returns max, $O(\log^2 n)$.
\texttt{querySubtree(v)} uses a half-open Euler interval, excluding $v$ in
edge mode. Initialize values: LazySegmentTree's default is $-10^9$, not zero.
Changing max to sum requires changing both layers. Noncommutative path
aggregates require directional processing. Depends on LazySegmentTree and its
BumpAllocator include. This longer source is not printed here.

\section{Exact geometry and constructions}
\subsection{Point operations}
\texttt{a.cross(b,c)} is $(b-a)\times(c-a)$, positive left turn.
For coordinate magnitude $B$, difference cross products can reach $8B^2$:
bound before using \texttt{Point<ll>}. Widen \emph{before} subtraction/product.
Integer points do not support fractional unit vectors or intersections exactly.
Use squared distances for ordering, with overflow checked.
\sourceimport{geometry}{Point.h}

\subsection{Closed segment membership}
Includes endpoints and degenerate segments. With floating coordinates, this
exact equality test may need the source's distance-based alternative.
\sourceimport{geometry}{OnSegment.h}

\subsection{Segment intersection}
Returns zero points for none, one for a single intersection, two for an overlap.
Integer point output truncates a nonintegral proper intersection; intermediates
include triple-coordinate products. For existence only, use orientation signs
and on-segment cases without constructing coordinates. Compare signs rather
than multiplying large cross products.
\sourceimport{geometry}{SegmentIntersection.h}

\subsection{Convex hull}
Returns CCW vertices, removing intermediate collinear boundary points.
Time $O(n\log n)$. Check size before subsequent hull operations: empty and
all-collinear inputs are degenerate.
\sourceimport{geometry}{ConvexHull.h}

\subsection{Signed doubled area}
Requires nonempty boundary-ordered polygon. Positive for CCW, negative clockwise;
self-intersections give algebraic area. Accumulation needs a bound beyond one
cross product. Keep doubled area integral for half-unit output.
\sourceimport{geometry}{PolygonArea.h}

\subsection{More geometry source contracts}
\textbf{InsidePolygon.h:} \texttt{inPolygon(poly,p,strict=true)}, $O(n)$,
simple polygon; false on boundary unless \texttt{strict=false}.
\textbf{PointInsideHull.h:} \texttt{inHull(h,p,strict)}, $O(\log n)$,
nonempty CCW convex hull with intermediate collinear points removed.
\textbf{HullDiameter.h:} nonempty hull, same order/collinearity conventions,
$O(n)$ after hull. \textbf{ClosestPair.h:} at least two points, $O(n\log n)$.
These are source references, not printed code. Exact predicates and approximate
coordinate constructions have different requirements; no universal EPS fixes both.

\section{Offline queries and optimized DP}
\subsection{Sweep contracts}
Sort points/queries by one coordinate; Fenwick aggregates the other compressed
coordinate. For $x\le X$, insert equal-$x$ points before queries; for $x<X$,
query first. Strict predecessor DP often needs batched equal-coordinate updates.
Compression preserves ordering, not actual lengths.

\textbf{FenwickTree2d.h:} call \texttt{fakeUpdate(x,y)} for all possible update
locations, then \texttt{init()} once. \texttt{query(X,Y)} sums $x<X,y<Y$;
rectangle sums use four prefixes. Operations $O(\log^2 n)$.
\textbf{MoQueries.h:} contains intentional callback placeholders.
Implement \texttt{add,del,calc}; intervals $[l,r)$ start empty.
The $O(n\sqrt q)$ movement estimate assumes cheap callbacks and a suitable
block size; literal 350 is not universal. Both sources are mapped, not printed.

\subsection{Maximum of lines}
Use only after deriving $dp_i=\max_j(k_jx_i+b_j)+c_i$.
\texttt{LineContainer lc; lc.add(k,b); lc.query(x)} returns a maximum;
empty queries assert. Negate slopes/intercepts/results for minima with overflow
checked. Arbitrary insertion/query order is allowed; time $O(\log n)$ amortized.
Eligibility still controls whether to insert before or after querying.
Products and internal breakpoint differences must fit \texttt{ll}.
\sourceimport{data-structures}{LineContainer.h}

\subsection{Small-state matrix powers}
Zero-initialized \texttt{Matrix<T,N> A}; fill \texttt{A.d}.
\texttt{(A\^{}k)*state} applies $k$ transitions to a column vector:
rows are destinations, columns sources. Time $O(N^3\log k)$.
Plain \texttt{ll} does not reduce modulo anything: use a modular numeric type
or adapt arithmetic. Exponent zero returns identity.
\sourceimport{data-structures}{Matrix.h}

\subsection{Transform and recurrence source contracts}
\textbf{FastFourierTransform.h:} \texttt{conv(vd,vd)} returns coefficient
convolution, $O(N\log N)$ at padded power-of-two $N$.
Sufficient rounded-integer precision bound:
$(\sum a_i^2+\sum b_i^2)\log_2N<9\cdot10^{14}$.
Round explicitly; do not promise exactness beyond the bound.
\textbf{NumberTheoreticTransform.h:} exact modulo 998244353, inputs normalized,
actual padded $N\le2^{23}$. This source pads \emph{strictly above} result size.
Modular output is not exact integer output without a coefficient bound or CRT.
Names such as \texttt{conv,mod} collide when mixing variants.
\textbf{LinearRecurrence.h:} \texttt{linearRec(S,tr,k)} computes $S_k$ for
$S_i=\sum_jS_{i-j-1}tr_j$, $O(d^2\log k)$. Source declares a mod-5 testing
placeholder marked \texttt{exclude-line}; supply the intended modulus when
copying. These longer sources are not printed here.
